% !TeX root = ../../Book5.tex
% 纲要标题：### 20.1 范畴 O
% 纲要位置：工作记录/计划纲要.md，第 4362 行。

\section{范畴 O}
\label{b5:sec:2001}

% 纲要范围：
%
% * 有限生成性、权分解及局部有限性
% * Verma 模和简单模
% * 本节需要第四卷的交换范畴语言
%

无限维最高权模仍有有限维权空间，也仍能用有限个向量生成。研究它们的扩张时，必须把这两个条件与正根方向的有限性同时保留。

\subsection{有限生成性、权分解与局部有限性}
\label{b5:subsec:2001:01}

\begin{definition}\label{b5:def:category-o}
设 \(\mathfrak g\) 为有限维复半单李代数，固定三角分解
\(\mathfrak g=\mathfrak n^-\oplus\mathfrak h\oplus\mathfrak n^+\)，其中 \(\mathfrak h\) 为 Cartan 子代数，\(\mathfrak n^\pm\) 分别为所选正、负根空间之和。在幺性左 \(U(\mathfrak g)\) 模范畴中，选取满足下列条件的模 \(M\)：
\begin{enumerate}
\item 模 \(M\) 作为 \(U(\mathfrak g)\) 模有限生成；
\item \(M=\bigoplus_{\mu\in\mathfrak h^*}M_\mu\)，其中 \(M_\mu\) 是共同特征值为 \(\mu\) 的 Cartan 权空间；
\item 对每个 \(v\in M\)，向量空间 \(U(\mathfrak n^+)v\) 有限维。
\end{enumerate}
这些模及它们之间的全部模同态组成的全子范畴，称为\term[fan-chou-O]{范畴 \(\mathcal O\)}[category O]。第三条称 \(\mathfrak n^+\) 的作用\term[jubu-youxian-zuoyong]{局部有限}[locally finite]。
\end{definition}

\begin{proposition}\label{b5:prop:o-weight-finiteness}
设 \(\mathfrak g\) 为有限维复半单李代数，固定 Cartan 子代数 \(\mathfrak h\) 及简单根 \(\alpha_1,\ldots,\alpha_r\)，令 \(\mathcal O\) 为相应的范畴 O，\(Q_+=\sum_i\mathbb Z_{\ge0}\alpha_i\)。设 \(M\in\mathcal O\)，记 \(M_\mu\) 为权 \(\mu\in\mathfrak h^*\) 的空间，\(\operatorname{wt}M\) 为全部非零权空间的权集合。则每个 \(M_\mu\) 有限维，且存在有限集 \(F\subset\mathfrak h^*\)，使
\[
\operatorname{wt}M\subseteq\bigcup_{\nu\in F}(\nu-Q_+).
\]
若 \(M\ne0\)，它含有一个非零最高权向量。
\end{proposition}
\begin{proof}
把一组有限生成元分别拆成有限个权分量，得到有限个权生成元 \(v_j\)。每个 \(U(\mathfrak n^+)v_j\) 有限维，而且由权向量生成，故被 \(\mathfrak h\) 保持。因此
\(E=\sum_jU(\mathfrak b)v_j=\sum_jU(\mathfrak n^+)v_j\)
是有限维 \(\mathfrak b\) 子模。PBW 给出 \(M=U(\mathfrak n^-)E\)。取 \(E\) 的权集为 \(F\)，即得权集包含式。

对固定 \(\mu\)，每个 \(\nu\in F\) 所贡献的负根单项式满足
\(\sum_{\alpha>0}a_\alpha\alpha=\nu-\mu\)。右端若不在 \(Q_+\)，没有贡献；若在 \(Q_+\)，所有 \(a_\alpha\) 之和至多为其高度，故只有有限种可能。\(E_\nu\) 也有限维，所以 \(M_\mu\) 有限维。

任取实际出现的权 \(\mu\)。处在 \(\mu+Q_+\) 中的权只有有限多个：对每个 \(\nu\in F\)，若 \(\mu+\beta=\nu-\gamma\)，则 \(\beta,\gamma\in Q_+\) 且 \(\beta+\gamma=\nu-\mu\)，高度受到固定上界约束。取这些权中的极大者 \(\eta\)，则全部正根向量都零化 \(M_\eta\)。任意非零 \(v\in M_\eta\) 就是所需向量。
\end{proof}

局部有限不表示整个模有限维。对 \(\mathfrak{sl}_2\)，Verma 模的 \(f\) 方向无限延伸，而 \(e\) 反复作用于任一向量最终为零。另一方面，去掉有限生成条件后，\(\bigoplus_{j\ge0}\mathbb C\) 虽有权分解且正根作用为零，零权空间却无限维。

\subsection{Verma 模与简单模}
\label{b5:subsec:2001:02}

\begin{proposition}\label{b5:prop:o-simple-objects}
设 \(\mathfrak g\) 为有限维复半单李代数，固定 Cartan 子代数 \(\mathfrak h\) 和正根集合，令 \(\mathcal O\) 为相应的范畴 O。对 \(\lambda\in\mathfrak h^*\)，记 \(M(\lambda)\) 为最高权 \(\lambda\) 的 Verma 模，\(L(\lambda)\) 为其唯一简单商。则这两个模都属于 \(\mathcal O\)，而 \(\mathcal O\) 的所有简单对象恰为这些 \(L(\lambda)\)，参数 \(\lambda\) 不重复。
\end{proposition}
\begin{proof}
由\cref{b5:prop:verma-pbw-universal}，\(M(\lambda)\) 是循环权模，权在 \(\lambda-Q_+\) 中且权空间有限维。对权为 \(\lambda-\beta\) 的向量，正根方向只能升到 \(\lambda-\beta+\gamma\le\lambda\)，其中 \(\gamma\in Q_+\) 的高度至多为 \(\beta\) 的高度；只有有限个这样的权。因此其 \(U(\mathfrak n^+)\) 轨道有限维。一般向量只有有限个权分量，结论仍成立。商模 \(L(\lambda)\) 继承三个条件。

若 \(S\) 是 \(\mathcal O\) 中的简单对象，\cref{b5:prop:o-weight-finiteness}给出最高权向量 \(v\)。子模 \(U(\mathfrak g)v\) 本身是最高权模，满足 \(\mathcal O\) 的条件，所以简单性使它等于 \(S\)。于是 Verma 泛性质和\cref{b5:thm:verma-simple-quotient}给出 \(S\cong L(\lambda)\)，并保证参数唯一。
\end{proof}

这里的分类允许任意复最高权。有限维简单模只占其中的支配整权部分。对于 \(\mathfrak{sl}_2\)，\eqref{b5:eq:sl2-verma-actions}说明：当 \(z\notin\mathbb Z_{\ge0}\) 时，\(M(z)\) 已经简单；当 \(z=m\ge0\) 为整数时，\(f^{m+1}v\) 生成它的唯一非零真子模。这一变化不能由有限维完全可约性解释，因为 Verma 模从来不是有限维模。

\subsection{交换范畴基本概念回顾}
\label{b5:subsec:2001:03}

一个范畴中的核与余核必须仍处在该范畴内，才可以在其中讨论正合列。下面检查这一点，同时指出 \(\mathcal O\) 与全部包络代数模之间的一处区别。

\begin{lemma}\label{b5:lem:enveloping-noetherian}
若 \(\mathfrak g\) 是有限维复李代数，则 \(U(\mathfrak g)\) 左、右 Noether。因此有限生成 \(U(\mathfrak g)\) 模的每个子模仍有限生成。
\end{lemma}
\begin{proof}
PBW 给出 \(\operatorname{gr}U(\mathfrak g)\cong\operatorname{Sym}(\mathfrak g)\)，后者为有限个变量的多项式环，由 Hilbert 基定理 Noether。对左理想 \(I\) 取诱导滤过，其主符号构成分次理想 \(\operatorname{gr}I\)。选择有限个齐次生成元，并在 \(I\) 中提升为 \(u_1,\ldots,u_s\)。任取 \(u\in I\)，先用 \(u_i\) 的左倍数消去最高次主符号，所得差次数严格下降；按次数归纳可把 \(u\) 写成这些 \(u_i\) 的左线性组合。因此左理想有限生成。右侧同理。最后对有限秩自由模的秩归纳，再把一般有限生成模表示为自由模的商，即得子模结论。
\end{proof}

\begin{proposition}\label{b5:prop:o-abelian}
设 \(\mathfrak g\) 为有限维复半单李代数，固定 Cartan 子代数及正根集合，令 \(\mathcal O\) 为相应的范畴 O。则 \(\mathcal O\) 是交换范畴，核、余核、有限直和以及像均按通常的 \(U(\mathfrak g)\) 模构造。
\end{proposition}
\begin{proof}
若 \(f:M\to N\) 是模同态，核与商继承权分解：对一个向量的有限个权分量，选 \(h\in\mathfrak h\) 区分这些权，再用 \(h\) 的插值多项式提取各分量，故子模也按权分解。局部 \(\mathfrak n^+\) 有限性对子模和商显然保持；有限生成性对商保持，对核由\cref{b5:lem:enveloping-noetherian}保持。有限直和也满足三个条件。于是模范畴中的核、余核都在 \(\mathcal O\) 内，且余像到像的自然映射仍为同构，这正是交换范畴的判据。
\end{proof}

不能据此断言它在全部 \(U(\mathfrak g)\) 模中对任意扩张封闭。例如令 \(\mathfrak g=\mathfrak{sl}_2\)，在二维 \(\mathfrak b\) 模 \(E\) 上令 \(e=0\)、\(h=\lambda I+N\)，其中 \(N\ne0\)、\(N^2=0\)。诱导模 \(U(\mathfrak g)\otimes_{U(\mathfrak b)}E\) 是两个 \(M(\lambda)\) 的扩张；诱导正合性来自 PBW 的右自由性。但其中 \(h\) 不是半单算子，所以中间项不在 \(\mathcal O\)。讨论 \(\operatorname{Ext}_{\mathcal O}\) 时，允许的中间对象也必须属于 \(\mathcal O\)。
